\documentclass[aspectratio=169]{beamer}

% -------------------------------------------------------
%  Paquetes
% -------------------------------------------------------
\usepackage[spanish]{babel}
\usepackage[utf8]{inputenc}
\usepackage{amsmath, amssymb, amsthm}
\usepackage{mathtools}
\usepackage{tikz}
\usetikzlibrary{arrows.meta, positioning, shapes.geometric, fit, backgrounds}
\usetikzlibrary{arrows.meta, positioning, shapes.geometric, fit, backgrounds, decorations.pathreplacing}
\usepackage{booktabs}
\usepackage{xcolor}


% -------------------------------------------------------
%  Tema y colores
% -------------------------------------------------------
\usetheme{Madrid}
\usecolortheme{seahorse}

\definecolor{azuloscuro}{RGB}{0,51,102}
\definecolor{azulmedio}{RGB}{0,102,179}
\definecolor{grisclaro}{RGB}{240,240,245}
\definecolor{rojoacento}{RGB}{180,30,30}
\definecolor{verdeacento}{RGB}{0,120,60}
\definecolor{act1}{RGB}{0,70,140}
\definecolor{act2}{RGB}{0,120,60}
\definecolor{puente}{RGB}{150,80,0}
\definecolor{act3}{RGB}{140,0,50}

\setbeamercolor{frametitle}{bg=azuloscuro, fg=white}
\setbeamercolor{title}{fg=white}
\setbeamercolor{author}{fg=white}
\setbeamercolor{institute}{fg=white}
\setbeamercolor{date}{fg=white}
\setbeamercolor{titlelike}{bg=azuloscuro}
\setbeamercolor{block title}{bg=azulmedio, fg=white}
\setbeamercolor{block body}{bg=grisclaro}
\setbeamercolor{alerted text}{fg=rojoacento}

% Bloques de colores para cada acto
\newenvironment{act1block}[1]{%
  \setbeamercolor{block title}{bg=act1, fg=white}%
  \setbeamercolor{block body}{bg=act1!10}%
  \begin{block}{#1}}{\end{block}}

\newenvironment{act2block}[1]{%
  \setbeamercolor{block title}{bg=act2, fg=white}%
  \setbeamercolor{block body}{bg=act2!10}%
  \begin{block}{#1}}{\end{block}}

\newenvironment{puenteblock}[1]{%
  \setbeamercolor{block title}{bg=puente, fg=white}%
  \setbeamercolor{block body}{bg=puente!10}%
  \begin{block}{#1}}{\end{block}}

\newenvironment{act3block}[1]{%
  \setbeamercolor{block title}{bg=act3, fg=white}%
  \setbeamercolor{block body}{bg=act3!10}%
  \begin{block}{#1}}{\end{block}}

% -------------------------------------------------------
%  Comandos matemáticos
% -------------------------------------------------------
\newcommand{\Prob}{\mathbb{P}}
\newcommand{\Exp}{\mathbb{E}}
\newcommand{\Geom}{\mathrm{Geom}}
\newcommand{\Expo}{\mathrm{Exp}}
\newcommand{\dconv}{\xrightarrow{d}}

% -------------------------------------------------------
%  Metadatos
% -------------------------------------------------------
\title[Tiempo de Primer Encuentro]{%
  Teoremas Límite para el Tiempo de\\
  Primer Encuentro del Óptimo Global}
\subtitle{Del modelo de Markov a la ley exponencial}
\author[]{Basado en Rudolph (1998) y Oliveto \& Witt (2015)}
\institute[]{Algoritmos Evolutivos --- Teoría}
\date{}

% -------------------------------------------------------
\begin{document}
% -------------------------------------------------------

% ==================== PORTADA ==========================
\begin{frame}
  \titlepage
\end{frame}

% ==================== ÍNDICE ===========================
\begin{frame}{Estructura de la exposición}
  \begin{center}
  \begin{tikzpicture}[
      box/.style={rectangle, rounded corners=3pt,
                  draw, fill=white, text width=2.4cm,
                  align=center, font=\scriptsize\bfseries,
                  inner sep=4pt, minimum height=1.0cm},
      arr/.style={-Stealth, thick, shorten >=2pt, shorten <=2pt},
    ]
    \node[box, draw=act1, fill=act1!15] (a1)
      {\textcolor{act1}{Acto 1}\\[2pt] El EA como\\Cadena de Markov};

    \node[box, draw=act2, fill=act2!15, right=0.7cm of a1] (a2)
      {\textcolor{act2}{Acto 2}\\[2pt] $\Prob\{T<\infty\}=1$\\bajo A$_1$--A$_3$};

    \node[box, draw=puente, fill=puente!15, right=0.7cm of a2] (pb)
      {\textcolor{puente}{Puente}\\[2pt] Cola geométrica\\de $T$};

    \node[box, draw=act3, fill=act3!15, right=0.7cm of pb] (a3)
      {\textcolor{act3}{Acto 3}\\[2pt] Mutación rara\\$\to$ Ley Exp$(1)$};

    \draw[arr, color=act1] (a1) -- (a2);
    \draw[arr, color=act2] (a2) -- (pb);
    \draw[arr, color=puente] (pb) -- (a3);
  \end{tikzpicture}
  \end{center}

  \bigskip
  \begin{columns}[T]
    \column{0.48\textwidth}
    \begin{block}{Referencias base}
      \begin{itemize}
        \item Rudolph (1998): marco general para EAs con
              espacio finito y tiempo discreto.
        \item Oliveto \& Witt (2015): análisis riguroso
              del SGA en OneMax.
      \end{itemize}
    \end{block}

    \column{0.48\textwidth}
    \begin{block}{Pregunta central}
      Dado un EA sobre un espacio finito,\\
      ¿cuál es la \textbf{distribución} del tiempo
      de espera hasta encontrar el óptimo?
    \end{block}
  \end{columns}
\end{frame}

% =======================================================
%  ACTO 1
% =======================================================
\section{Acto 1 — El EA como Cadena de Markov}

\begin{frame}{Acto 1 \textcolor{act1}{\rule[0.5ex]{6cm}{1.5pt}}}
  \begin{center}
    \Large\textcolor{act1}{\textbf{¿Cómo modelamos el problema?}}\\[8pt]
    \normalsize El algoritmo evolutivo como cadena de Markov finita
  \end{center}
\end{frame}

\begin{frame}{El EA como proceso estocástico}
  Un algoritmo evolutivo transforma la población en cada generación
  mediante operadores probabilísticos:

  \bigskip
  \begin{center}
  \begin{tikzpicture}[
      node distance=0.5cm,
      op/.style={rectangle, rounded corners=2pt,
                 draw=act1, fill=act1!10,
                 font=\tiny, inner sep=4pt, align=center,
                 minimum width=1.6cm},
      arr/.style={-Stealth, thick, color=act1, shorten >=1pt, shorten <=1pt}
    ]
    \node[op] (p)  {Población\\$X_k$};
    \node[op, right=0.6cm of p] (m)  {Selección\\padres};
    \node[op, right=0.6cm of m] (r)  {Recom-\\binación};
    \node[op, right=0.6cm of r] (mu) {Mutación};
    \node[op, right=0.6cm of mu] (s) {Selección\\hijos};
    \node[op, right=0.6cm of s] (p2) {Población\\$X_{k+1}$};

    \draw[arr] (p)  -- (m);
    \draw[arr] (m)  -- (r);
    \draw[arr] (r)  -- (mu);
    \draw[arr] (mu) -- (s);
    \draw[arr] (s)  -- (p2);
  \end{tikzpicture}
  \end{center}

  \bigskip
  \begin{act1block}{Observación clave (Rudolph, 1998)}
    $X_{k+1}$ depende \textbf{únicamente} de $X_k$, no de
    $X_0, X_1, \ldots, X_{k-1}$. Esto es la
    \textbf{propiedad de Markov}.
  \end{act1block}
\end{frame}

\begin{frame}{Cadena de Markov finita}
  \begin{act1block}{Modelo formal}
    \[
      \Prob(X_{k+1} \in A \mid X_0, X_1, \ldots, X_k)
      \;=\;
      \Prob(X_{k+1} \in A \mid X_k).
    \]
  \end{act1block}

  \bigskip
  Como el espacio de búsqueda $\mathcal{X}$ es \textbf{finito}, el
  espacio de estados de la cadena también lo es. Esto nos permite:

  \bigskip
  \begin{columns}[T]
    \column{0.48\textwidth}
    \textbf{Lo que ganamos con este modelo:}
    \begin{itemize}
      \item Usar toda la maquinaria de cadenas de Markov finitas.
      \item Definir rigurosamente el tiempo de primer encuentro.
      \item Estudiar la distribución del tiempo de espera.
    \end{itemize}

    \column{0.48\textwidth}
    \begin{act1block}{Tiempo de primer encuentro}
      \[
        T = \min\{k \geq 0 : F_k = f^*\}
      \]
      donde $F_k = \max\{f(X_k^{(i)})\}$ es el mejor
      fitness en la generación $k$.\\[4pt]
      $T$ es un \textbf{tiempo de paro} bien definido.
    \end{act1block}
  \end{columns}
\end{frame}

% =======================================================
%  ACTO 2
% =======================================================
\section{Acto 2 — Llegada al óptimo con probabilidad 1}

\begin{frame}{Acto 2 \textcolor{act2}{\rule[0.5ex]{6cm}{1.5pt}}}
  \begin{center}
    \Large\textcolor{act2}{\textbf{¿Qué garantizamos?}}\\[8pt]
    \normalsize Bajo qué suposiciones el EA encuentra el óptimo con $\Prob = 1$
  \end{center}
\end{frame}

\begin{frame}{Suposiciones sobre los operadores}
  Para garantizar la llegada al óptimo necesitamos condiciones mínimas
  sobre cada operador. Rudolph (1998) formula tres:

  \bigskip
  \begin{columns}[T]
    \column{0.32\textwidth}
    \begin{act2block}{(A$_1$) Recombinación}
      Todo padre puede pasar la recombinación \textbf{sin alterarse}
      con probabilidad mínima:
      \[\delta_r > 0\]
    \end{act2block}

    \column{0.32\textwidth}
    \begin{act2block}{(A$_2$) Mutación}
      Todo individuo puede alcanzar \textbf{cualquier otro} mediante
      mutaciones sucesivas con probabilidad mínima por paso:
      \[\delta_m > 0\]
    \end{act2block}

    \column{0.32\textwidth}
    \begin{act2block}{(A$_3$) Selección}
      Todo individuo puede \textbf{sobrevivir} la selección con
      probabilidad mínima:
      \[\delta_s > 0\]
    \end{act2block}
  \end{columns}

  \bigskip
  \begin{itemize}
    \item Ninguna suposición puede ser eliminada: si alguna
          $\delta_r, \delta_m$ o $\delta_s$ es \textbf{cero},
          el resultado falla.
  \end{itemize}
\end{frame}

\begin{frame}{¿Por qué cada suposición es necesaria?}
  \begin{center}
  \begin{tabular}{clll}
    \toprule
    \textbf{Si faltara\ldots} & \textbf{Consecuencia} \\
    \midrule
    (A$_1$): $\delta_r = 0$ &
      La recombinación destruye siempre al individuo; nunca avanza. \\[4pt]
    (A$_2$): $\delta_m = 0$ &
      El óptimo puede ser inalcanzable desde algún estado inicial. \\[4pt]
    (A$_3$): $\delta_s = 0$ &
      Un hijo que llegó al óptimo siempre es eliminado por selección. \\
    \bottomrule
  \end{tabular}
  \end{center}

  \bigskip
  \begin{act2block}{Intuición}
    Las tres suposiciones juntas garantizan que el espacio de búsqueda
    está \textbf{completamente conectado} a través de los operadores,
    y que ningún operador destruye determinísticamente el progreso.
  \end{act2block}
\end{frame}

\begin{frame}{La prueba del Teorema 2.1 (Rudolph, 1998)}
  La prueba construye el argumento en tres niveles:

  \medskip
  \begin{enumerate}
    \item<1-> \textbf{Existe un camino hacia el óptimo.}\\
          Por (A$_2$), para cualquier $x \notin \mathcal{X}^*$ existe
          una secuencia finita de mutaciones:
          \[x = x_0 \to x_1 \to \cdots \to x_{k^*} \in \mathcal{X}^*\]
          Como $\mathcal{X}$ es finito, $k^* = \max_x\{k_x\}$ es finito.

    \bigskip
    \item<2-> \textbf{La probabilidad de recorrer ese camino es positiva.}\\
          En cada paso: recombinación no altera ($\geq \delta_r$),
          mutación avanza ($\geq \delta_m$), selección deja sobrevivir
          ($\geq \delta_s$). En $k^*$ pasos:
          \[\delta = (\delta_r \cdot \delta_m \cdot \delta_s)^{k^*-1}
            \cdot \delta_r \cdot \delta_m > 0\]

    \bigskip
    \item<3-> \textbf{Con suficientes intentos, el éxito es inevitable.}\\
          En cada bloque de $k^*$ iteraciones, $\Prob(\text{éxito}) \geq \delta > 0$,
          independientemente del estado actual (propiedad de Markov).
  \end{enumerate}
\end{frame}

\begin{frame}{Teorema 2.1 (Rudolph, 1998)}
  \begin{theorem}[Rudolph, 1998]
    Si se cumplen \emph{(A$_1$)}, \emph{(A$_2$)} y \emph{(A$_3$)},
    entonces:
    \[
      \Prob\{T < \infty\} = 1
    \]
    independientemente de la inicialización. Además, la probabilidad
    de \textbf{no} haber encontrado el óptimo después de $k$
    iteraciones satisface:
    \[
      \Prob\{T > k\}
      \;\leq\;
      (1-\delta)^{\lfloor k/k^* \rfloor}
      \;\xrightarrow{k\to\infty}\; 0.
    \]
  \end{theorem}

  \bigskip
  \begin{columns}[T]
    \column{0.48\textwidth}
    \textbf{Lo que ya sabemos:}
    \begin{itemize}
      \item El óptimo se encuentra en tiempo finito, siempre.
      \item La probabilidad de no haberlo encontrado cae a cero.
    \end{itemize}

    \column{0.48\textwidth}
    \textbf{Lo que aún no sabemos:}
    \begin{itemize}
      \item ¿Qué \textit{forma} tiene la distribución de $T$?
      \item ¿Cómo se comporta cuando $\delta_m \to 0$?
    \end{itemize}
  \end{columns}
\end{frame}

% =======================================================
%  PUENTE
% =======================================================
\section{Puente — Cola geométrica de T}

\begin{frame}{Puente \textcolor{puente}{\rule[0.5ex]{6cm}{1.5pt}}}
  \begin{center}
    \Large\textcolor{puente}{\textbf{La cola geométrica de $T$}}\\[8pt]
    \normalsize Llegar al infinito
  \end{center}
\end{frame}

\begin{frame}{La estructura de bloques}
  La prueba del Teorema 2.1 divide el tiempo en
  \textbf{bloques de $k^*$ iteraciones}:

  \bigskip
  \begin{center}
  \begin{tikzpicture}[
      bloque/.style={rectangle, draw=puente, fill=puente!10,
                     minimum width=2.2cm, minimum height=0.9cm,
                     font=\small, align=center},
      arr/.style={-Stealth, thick, color=puente}
    ]
    \node[bloque] (b1) {Bloque 1\\$k^*$ iter.};
    \node[bloque, right=0.3cm of b1] (b2) {Bloque 2\\$k^*$ iter.};
    \node[right=0.3cm of b2, font=\large] (dots) {$\cdots$};
    \node[bloque, right=0.3cm of dots] (bj) {Bloque $j$\\$k^*$ iter.};
    \node[right=0.3cm of bj, font=\large] (dots2) {$\cdots$};

    \draw[arr] (b1) -- (b2);
    \draw[arr] (b2) -- (dots);
    \draw[arr] (dots) -- (bj);
    \draw[arr] (bj) -- (dots2);

    \node[below=0.5cm of b1, font=\small, color=rojoacento]
      {$\Prob(\text{éxito})\geq\delta$};
    \node[below=0.5cm of b2, font=\small, color=rojoacento]
      {$\Prob(\text{éxito})\geq\delta$};
    \node[below=0.5cm of bj, font=\small, color=rojoacento]
      {$\Prob(\text{éxito})\geq\delta$};
  \end{tikzpicture}
  \end{center}

  \bigskip
  \begin{puenteblock}{Independencia entre bloques}
    Por la propiedad de Markov, la probabilidad de éxito $\delta$ en
    cada bloque es \textbf{la misma e independiente} del resultado de
    bloques anteriores, sin importar cuán cerca o lejos del óptimo
    se esté.
  \end{puenteblock}
\end{frame}

\begin{frame}{¿Por qué esto es una distribución geométrica?}
  Cada bloque es equivalente a \textbf{lanzar una moneda} con
  probabilidad de éxito $\delta$:

  \bigskip
  \begin{columns}[T]
    \column{0.48\textwidth}
    \textbf{Lanzamiento de moneda:}
    \begin{itemize}
      \item Éxito (prob. $\delta$): encontrar el óptimo en este bloque.
      \item Fracaso (prob. $1-\delta$): no encontrarlo, pasar al siguiente.
    \end{itemize}

    \bigskip
    Define $B = \lfloor T/k^* \rfloor$ como el número de bloques
    hasta el primer éxito. Entonces:
    \[
      \Prob\{B = j\} = (1-\delta)^{j-1} \cdot \delta
    \]
    que es una \textbf{distribución geométrica}.

    \column{0.48\textwidth}
    \begin{puenteblock}{$T$ hereda la estructura}
      Como $T \approx k^* \cdot B$:
      \[
        \Prob\{T > k\} \leq (1-\delta)^{\lfloor k/k^* \rfloor}
      \]
      \[
        \Exp[T] \approx \frac{k^*}{\delta}
      \]
      La cola de $T$ cae \textbf{exponencialmente} en $k$:
      \[
        (1-\delta)^k \approx e^{-\delta k}
      \]
    \end{puenteblock}
  \end{columns}
\end{frame}

\begin{frame}{¿Qué significa $\lfloor k/k^* \rfloor$?}
  Es el número de \textbf{bloques completos} que caben en $k$
  iteraciones. Las iteraciones sobrantes se ignoran porque no
  garantizan completar el camino al óptimo.

  \bigskip
  \textbf{Ejemplo con $k^* = 3$:}

  \begin{center}
  \begin{tikzpicture}[
      iter/.style={rectangle, draw=gray, fill=gray!10,
                   minimum width=0.7cm, minimum height=0.6cm,
                   font=\scriptsize},
      brace/.style={decorate,
                    decoration={brace, amplitude=5pt, mirror}}
    ]
    \foreach \i in {1,...,9} {
      \node[iter] at ({(\i-1)*0.85cm}, 0) {\i};
    }
    \node[iter, fill=orange!20, draw=orange] at ({6*0.85cm}, 0) {7};
    \node[iter, fill=gray!5, draw=gray!40] at ({7*0.85cm}, 0) {8};
    \node[iter, fill=gray!5, draw=gray!40] at ({8*0.85cm}, 0) {9};

    \draw[brace, color=act2]
      (-0.35cm, -0.45cm) -- (2.2cm, -0.45cm)
      node[midway, below=5pt, font=\scriptsize, color=act2]
      {Bloque 1};
    \draw[brace, color=act2]
      (2.2cm, -0.45cm) -- (4.75cm, -0.45cm)
      node[midway, below=5pt, font=\scriptsize, color=act2]
      {Bloque 2};
    \draw[brace, color=puente]
      (4.75cm, -0.45cm) -- (5.5cm, -0.45cm)
      node[midway, below=5pt, font=\scriptsize, color=puente]
      {sobrante};
  \end{tikzpicture}
  \end{center}

  \bigskip
  Para $k=7$ y $k^*=3$: $\lfloor 7/3 \rfloor = 2$ bloques completos.
  La iteración 7 cae en el tercer bloque incompleto y se ignora
  en la cota. Por eso:
  \[
    \Prob\{T > 7\} \leq (1-\delta)^2,
    \quad\text{no}\quad (1-\delta)^7.
  \]
\end{frame}

% =======================================================
%  ACTO 3
% =======================================================
\section{Acto 3 — Mutación rara y ley exponencial}

\begin{frame}{Acto 3 \textcolor{act3}{\rule[0.5ex]{6cm}{1.5pt}}}
  \begin{center}
    \Large\textcolor{act3}{\textbf{¿Qué pasa cuando la mutación es rara?}}\\[8pt]
    \normalsize De la distribución geométrica a la ley exponencial
  \end{center}
\end{frame}

\begin{frame}{¿Qué significa mutación rara?}
  \begin{columns}[T]
    \column{0.48\textwidth}
    \textbf{Definición intuitiva}\\[6pt]
    La mutación es \emph{rara} cuando la probabilidad de que un bit
    específico cambie en una iteración es \textbf{muy pequeña}: el
    individuo pasa casi sin cambios a la siguiente generación.

    \bigskip
    \textbf{En Rudolph:}
    \[\delta_m > 0, \quad \delta_m \approx 0.\]

    \textbf{En Oliveto \& Witt (Standard Bit Mutation):}
    \[p = \frac{1}{n} \;\xrightarrow{n\to\infty}\; 0.\]

    \column{0.48\textwidth}
    \begin{act3block}{¿Por qué $p=1/n$ es ``rara''?}
      Número esperado de bits que cambian:
      \[
        \Exp[\text{bits mutados}] = n \cdot \frac{1}{n} = 1.
      \]
      Para $n$ grande, \textbf{solo 1 bit de $n$} cambia en promedio
      por iteración.
    \end{act3block}

    \vspace{0.3cm}
    \begin{alertblock}{Consecuencia sobre $\delta$}
      \small
      $\delta_m \to 0 \;\Rightarrow\; \delta \to 0 \;\Rightarrow\;
      \Exp[T] = k^*/\delta \to \infty$
    \end{alertblock}
  \end{columns}
\end{frame}

\begin{frame}{De la geométrica a la exponencial}
  Cuando $\delta \to 0$, la distribución geométrica de $T$
  converge a una ley exponencial. Esto es un resultado clásico
  de probabilidad:

  \bigskip
  \begin{act3block}{Resultado clásico}
    Si $T \sim \Geom(\delta)$ con $\delta \to 0$, entonces:
    \[
      \delta \cdot T \;\dconv\; \Expo(1).
    \]
    Es decir, el tiempo reescalado por su media $\Exp[T]=1/\delta$
    converge en distribución a una exponencial estándar.
  \end{act3block}

  \bigskip
  \textbf{Demostración:}
  \[
    \Prob\{\delta \cdot T > x\}
    = \Prob\!\left\{T > \frac{x}{\delta}\right\}
    \leq (1-\delta)^{x/\delta}
    = \left[(1-\delta)^{1/\delta}\right]^x
    \;\xrightarrow{\delta\to 0}\; e^{-x}.
  \]
  Esto es la función de supervivencia de $\Expo(1)$.
\end{frame}



% =======================================================
%  SÍNTESIS
% =======================================================
\section{Síntesis}

\begin{frame}{La cadena lógica completa}
  \begin{center}
  \begin{tikzpicture}[
      node distance=0.5cm,
      box/.style={rectangle, rounded corners=4pt,
                  text width=10cm, align=center,
                  font=\small, inner sep=5pt,
                  minimum height=0.75cm},
      arr/.style={-Stealth, thick},
    ]
    \node[box, draw=act1, fill=act1!10] (m)
      {Propiedad de Markov $\;\Rightarrow\;$
       $T$ bien definido como tiempo de paro};

    \node[box, draw=act2, fill=act2!10, below=of m] (t)
      {Teorema 2.1 (Rudolph) $\;\Rightarrow\;$
       $\Prob\{T<\infty\}=1$ bajo (A$_1$), (A$_2$), (A$_3$)};

    \node[box, draw=puente, fill=puente!10, below=of t] (g)
      {Cola geométrica $\;\Rightarrow\;$
       $\Prob\{T>k\} \leq (1-\delta)^{\lfloor k/k^*\rfloor}$};

    \node[box, draw=act3, fill=act3!10, below=of g] (r)
      {Mutación rara $(\delta_m\to 0)$ $\;\Rightarrow\;$
       $\delta\to 0,\quad \mathbb{E}[T]\to\infty$};

    \node[box, draw=act3, fill=act3!15, below=of r] (e)
      {$\Geom(\delta)$ con $\delta\to 0$
       $\;\Rightarrow\;$
       $\dfrac{T}{\mathbb{E}[T]} \xrightarrow{d} \mathrm{Exp}(1)$};



    \draw[arr, color=act1]   (m)  -- (t);
    \draw[arr, color=act2]   (t)  -- (g);
    \draw[arr, color=puente] (g)  -- (r);
    \draw[arr, color=act3]   (r)  -- (e);

  \end{tikzpicture}
  \end{center}
\end{frame}

\begin{frame}{Conclusiones}
  \begin{enumerate}
    \item \textcolor{act1}{\textbf{Acto 1.}}
          La propiedad de Markov permite modelar cualquier EA con
          espacio finito como una cadena de Markov, haciendo que
          $T$ sea un tiempo de paro bien definido.

    \bigskip
    \item \textcolor{act2}{\textbf{Acto 2.}}
          Bajo las condiciones mínimas (A$_1$)--(A$_3$), el
          Teorema 2.1 de Rudolph garantiza $\Prob\{T<\infty\}=1$
          para cualquier inicialización.

    \bigskip
    \item \textcolor{puente}{\textbf{Puente.}}
          La prueba produce además una cota geométrica sobre la cola
          de $T$, que es el eslabón formal hacia la ley exponencial.

    \bigskip
    \item \textcolor{act3}{\textbf{Acto 3.}}
          Bajo mutación rara ($\delta_m \to 0$), el parámetro
          $\delta \to 0$ y el tiempo de primer encuentro reescalado
          converge en distribución a una \alert{ley exponencial}:
          \[
            \frac{T}{\mathbb{E}[T]}
            \;\xrightarrow{d}\;
            \mathrm{Exp}(1).
          \]
  \end{enumerate}
\end{frame}

\begin{frame}{Referencias}
  \begin{thebibliography}{9}
    \bibitem{rudolph1998}
      G.~Rudolph.
      \newblock Finite Markov Chain Results in Evolutionary Computation:
        A Tour d'Horizon.
      \newblock \emph{Fundamenta Informaticae}, in press, 1998.

    \bibitem{oliveto2015}
      P.\,S.~Oliveto, C.~Witt.
      \newblock Improved Time Complexity Analysis of the Simple Genetic
        Algorithm.
      \newblock \emph{Theoretical Computer Science}, 605:21--41, 2015.

    \bibitem{feller1971}
      W.~Feller.
      \newblock \emph{An Introduction to Probability Theory and Its
        Applications}, vol.~2.
      \newblock Wiley, 1971.
  \end{thebibliography}
\end{frame}

\begin{frame}{}
  \begin{center}
    {\Large \textbf{Gracias}}\\[12pt]
    \large Por favor no hagan preguntas ;) 
  \end{center}
\end{frame}

\end{document}